%----------------------------------------------------------------------------------------
% NOTATIE EN FIGUREN VAN DE CURSUS MATRICES
% Staat in de preambule en niet in een hoofdstuk: ook de lijsten met oplossingen gebruiken dit.
%----------------------------------------------------------------------------------------
% Bewaarde grafiek, geplaatst door maak-pdf.py: bestand en pagina in figuren-cache/
\newcommand{\grafiekbestand}[2]{\includegraphics[page=#2,trim=3bp 3bp 3bp 3bp]{figuren-cache/#1}}
\usetikzlibrary{matrix,positioning,backgrounds}

% Getransponeerde: A\tr
\newcommand{\tr}{^{\mkern-1.5mu\mathsf{T}}}
% Compacte matrix in lopende tekst en in de antwoorden: \km{1 & 2 \\ 3 & 4}
% Getallen op gewone grootte, maar met minder ruimte tussen rijen en kolommen dan bmatrix.
\newcommand{\km}[1]{{\setlength{\arraycolsep}{2.4pt}\renewcommand{\arraystretch}{.82}\left[\begin{array}{*{10}{c}}#1\end{array}\right]}}
% Matrix met wat meer ruimte tussen de kolommen
\newenvironment{mat}{\setlength{\arraycolsep}{4pt}\begin{bmatrix}}{\end{bmatrix}}

% Zoals stappen, maar met een bredere kolom voor de berekening: voor lange formules
\newenvironment{stappenbreed}{%
  \par\renewcommand{\arraystretch}{1.25}%
  \noindent\begin{tabular}{@{}>{$\displaystyle}p{9.8cm}<{$}@{\hspace{4mm}}>{\RaggedRight\small\itshape\arraybackslash}p{4.9cm}@{}}%
}{\end{tabular}\par}

% Zoals stappen, met eigen kolombreedtes: \begin{stappenvrij}{6.2cm}{3.8cm}. Voor naast een figuur.
\newenvironment{stappenvrij}[2]{%
  \par\renewcommand{\arraystretch}{1.25}%
  \noindent\begin{tabular}{@{}>{$\displaystyle}p{#1}<{$}@{\hspace{3mm}}>{\RaggedRight\small\itshape\arraybackslash}p{#2}@{}}%
}{\end{tabular}\par}

% Tekst of berekening links, figuur rechts: \naast[breedte rechts]{links}{rechts}
% De breedte rechts komt uit de figuur zelf; het optionele argument wordt niet gebruikt.
\newsavebox{\naastbox}
\newcommand{\naast}[3][]{%
  \par\noindent\sbox{\naastbox}{#3}%
  \begin{minipage}[t]{\dimexpr\linewidth-\wd\naastbox-3mm\relax}\vspace{0pt}#2\end{minipage}\hfill
  \begin{minipage}[t]{\wd\naastbox}\vspace{0pt}\usebox{\naastbox}\end{minipage}\par\addvspace{1mm}}
% Een rij figuren, gecentreerd
\newcommand{\figuren}[1]{\par\addvspace{1mm}{\centering #1\par}\addvspace{1.5mm}}
% Opgave met een grafiek rechts: \begin{metgrafiek}[breedte]{tikz} opgave en deelvragen \end{metgrafiek}
% De solution komt na \end{metgrafiek}, zoals bij metprent.
\newsavebox{\grafiekbox}
\newenvironment{metgrafiek}[2][]{%
  \sbox{\grafiekbox}{#2}%
  \begin{minipage}[t]{\dimexpr\linewidth-\wd\grafiekbox-3mm\relax}%
}{%
  \end{minipage}\hfill
  \begin{minipage}[t]{\wd\grafiekbox}\vspace{-2mm}\usebox{\grafiekbox}\end{minipage}\par}

% Rooster met assen: \rooster{xmin}{xstap}{xmax}{ymin}{ystap}{ymax}, met getallen op de assen (gehele waarden).
% \roosterleeg doet hetzelfde zonder getallen, voor eigen labels.
\newcommand{\roosterleeg}[6]{%
  \draw[black!14,very thin,xstep=#2,ystep=#5] (#1,#4) grid (#3,#6);
  \draw[->,grijs] (#1,0) -- ({#3+.4*#2},0) node[right,text=black] {$x$};
  \draw[->,grijs] (0,#4) -- (0,{#6+.4*#5}) node[above,text=black] {$y$};}
\newcommand{\rooster}[6]{%
  \roosterleeg{#1}{#2}{#3}{#4}{#5}{#6}
  \pgfmathtruncatemacro{\rtweex}{#1+#2}\pgfmathtruncatemacro{\rtweey}{#4+#5}
  \foreach \t in {#1,\rtweex,...,#3}{\ifnum\t=0\else\node[below,inner sep=1.5pt,text=black] at (\t,0) {$\t$};\fi}
  \foreach \t in {#4,\rtweey,...,#6}{\ifnum\t=0\else\node[left,inner sep=1.5pt,text=black] at (0,\t) {$\t$};\fi}}

% GRAFEN. Knopen en bogen in de kleur van het hoofdstuk.
%   \node[knoop] (a) at (0,0) {A};   \draw[boog] (a) -- (b);   \draw[boog] (a) to[bend left=20] (b);
%   lus: \draw[boog] (a) to[lus] (a);   kans bij een boog: node[kans] {0,3}
\tikzset{
  knoop/.style={circle,draw=accent,fill=accent!12,thick,minimum size=7mm,inner sep=0pt,font=\small\bfseries},
  boog/.style={-{Stealth[length=2.2mm]},accent,thick,shorten >=1pt,shorten <=1pt},
  lijn/.style={accent,thick},
  lus/.style={out=#1+25,in=#1-25,looseness=7,min distance=7mm},
  lus/.default=90,
  kans/.style={fill=white,inner sep=1pt,font=\scriptsize,text=black},
  % Matrix van getallen in TikZ, voor figuren waarin een rij of kolom gemarkeerd wordt
  % kolommen scheiden met \& (ampersand replacement), want de figuren staan vaak in een macro-argument
  getallen/.style={matrix of math nodes,ampersand replacement=\&,left delimiter={[},right delimiter={]},
    nodes={minimum width=6mm,minimum height=4.6mm,inner sep=1pt,anchor=center},column sep=1mm,row sep=.6mm},
}
% Markeer in een TikZ-matrix (naam m) een rij of kolom: \markeer{m-1-1}{m-1-3}
\newcommand{\markeer}[3][fel!45]{\begin{scope}[on background layer]\fill[#1,rounded corners=1.5pt] ([shift={(-.6mm,-.6mm)}]#2.south west) rectangle ([shift={(.6mm,.6mm)}]#3.north east);\end{scope}}
% Staaf in een staafdiagram: \staaf{x}{hoogte}{breedte}
\newcommand{\staaf}[3]{\fill[accent!55] ({#1-#3/2},0) rectangle ({#1+#3/2},#2); \draw[accent] ({#1-#3/2},0) rectangle ({#1+#3/2},#2);}
